% Part: lambda-calculus
% Chapter: church-rosser
% Section: definitions-and-properties

\documentclass[../../../include/open-logic-section]{subfiles}

\begin{document}

\olfileid{lam}{cr}{dap}

\olsection{সংজ্ঞা ও ধর্মাবলি}

এই অধ্যায়ে আমরা চার্চ--রসার ধর্মের ধারণা এবং এই ধর্মের কয়েকটি সাধারণ
বৈশিষ্ট্য আলোচনা করব।

\begin{defn}[চার্চ--রসার ধর্ম, CR] পদগুলির উপর কোনো সম্বন্ধ~$\xredone$
  \emph{চার্চ--রসার ধর্ম} পূরণ করে ঠিক তখনই, যখন $M \xredone P$ ও
  $M \xredone Q$ হলেই এমন কোনো~$N$ থাকে যাতে $P \xredone N$ এবং
  $Q \xredone N$।
\end{defn}

ল্যাম্বডা ক্যালকুলাসকে গণনার এমন একটি মডেল হিসেবে দেখা যায়, যেখানে
স্বাভাবিক রূপের পদগুলি হলো “মান” এবং পদগুলির উপর কোনো হ্রাসযোগ্যতা-সম্বন্ধ
হলো “গণনার বিধি”। চার্চ--রসার ধর্ম বলে: কোনো গণনা এগিয়ে নেওয়ার একাধিক
পথ থাকলেও অভিব্যক্তিটির মান একটিই থাকে।

প্রাথমিক বীজগণিত থেকে উদাহরণ নিই। $4 \times (1+2) + 3$ গণনা করার
একাধিক পথ আছে। প্রথমে $1+2$-কে~$3$-এ হ্রাস করলে এটি $4 \times 3+3$-এ
হ্রাস পায়; আবার বণ্টন ধর্ম ব্যবহার করে প্রথমে $4 \times (1+2)$ হ্রাস
করলে $4 \times 1+4 \times 2+3$ পাওয়া যায়। কিন্তু উভয়কেই আরও হ্রাস
করে $12+3$ পাওয়া যায়।

$\xredone$-কে $\beta$-হ্রাস ধরলে সহজেই দেখা যায়, চার্চ--রসার ধর্মের
একটি ফল হলো: কোনো পদের স্বাভাবিক রূপ থাকলে সেটি একক। ধরি, $M$-কে
$P$ ও~$Q$-তে হ্রাস করা যায় এবং উভয়েই স্বাভাবিক রূপ। চার্চ--রসার ধর্ম
অনুসারে এমন কোনো~$N$ আছে যাতে $P$ ও~$Q$ উভয়কেই তাতে হ্রাস করা যায়।
অনুমান অনুযায়ী $P$ ও~$Q$ স্বাভাবিক রূপ, তাই $P$ ও~$Q$-কে~$N$-এ হ্রাস কেবল
তুচ্ছ হ্রাসই হতে পারে; অর্থাৎ $P$, $Q$ ও~$N$ অভিন্ন। এই কারণেই আমরা
কোনো পদের \emph{স্বাভাবিক রূপটি} বলতে পারি।

অতএব, ল্যাম্বডা ক্যালকুলাসকে গণনার মডেল হিসেবে দেখলে কোনো পদের স্বাভাবিক
রূপকে ওই পদ থেকে শুরু হওয়া গণনার “চূড়ান্ত ফল” মনে করা যায়। উপরের
অনুসিদ্ধান্তটি বলে, এমন চূড়ান্ত ফল থাকলে তা একটিই; ঠিক যেমন
$4 \times (1+2)+3$ গণনার ফল একটিই, যথা~$15$।

\begin{thm} \ollabel{thm:str}
  কোনো সম্বন্ধ~$\xredone$ চার্চ--রসার ধর্ম পূরণ করলে, এবং $\xred$ যদি
  $\xredone$-কে ধারণকারী ক্ষুদ্রতম পরিযায়ী সম্বন্ধ হয়, তবে $\xred$-ও
  চার্চ--রসার ধর্ম পূরণ করে।
\end{thm}

\begin{proof}
  ধরি $M \xred P$ এবং $M \xred Q$; অর্থাৎ
  \begin{align*}
    M & \xredone P_1 \xredone \dots \xredone P_m \, (=P) \text{ এবং}\\
    M & \xredone Q_1 \xredone \dots \xredone Q_n \, (=Q).
  \end{align*}
% BN-SRC-292 TARGET-CORRECTION-BEGIN
\par\smallskip\noindent\textnormal{\small\textbf{উৎস-সংশোধন (BN-SRC-292):} হিমায়িত প্রমাণটি $P_1,\dots,P_m$ ও~$Q_1,\dots,Q_n$ ধারাদুটি লিখে পরে হঠাৎ বলে $N_{m,0}$ হলো~$P$ এবং $N_{0,n}$ হলো~$Q$, কিন্তু $P$ ও~$Q$-কে ধারার প্রান্তবিন্দু হিসেবে আগে চিহ্নিত করে না। বাংলা প্রমাণে শুরুতেই নির্বিচার $M\xred P$ ও~$M\xred Q$ নেওয়া হয়েছে এবং $P_m=P$, $Q_n=Q$ স্পষ্ট করা হয়েছে; এটিই $\xred$-এর চার্চ--রসার ধর্ম প্রমাণের প্রয়োজনীয় নির্বিচার যুগল।} % BN-SRC-292 TARGET-CORRECTION-END
  $m + 1$ উচ্চতা ও~$n + 1$ প্রস্থের পদগুলির একটি ছক~$N$ নির্মাণ করে
  উপপাদ্যটি প্রমাণ করব। $i$-তম সারি ও~$j$-তম স্তম্ভের পদকে $N_{i,j}$
  লিখি।

  $N_{i,j} \xredone N_{i+1,j}$ এবং $N_{i,j} \xredone N_{i,j+1}$ হয়
  এমনভাবে $N$ নির্মাণ করি। সংজ্ঞাটি নিচের মতো:
  \begin{align*}
    N_{0,0} &= M \\
    N_{i,0} &= P_i && \text{যদি } 1 \le i \le m \\
    N_{0,j} &= Q_j && \text{যদি } 1 \le j \le n \\
  \intertext{এবং অন্য ক্ষেত্রে:}
    N_{i,j} &= R
  \end{align*}
  যেখানে $R$ এমন একটি পদ যাতে $N_{i-1,j} \xredone R$ এবং $N_{i,j-1}
  \xredone R$। $\xredone$-এর চার্চ--রসার ধর্ম অনুসারে এমন পদ সর্বদা
  থাকে।

  এখন $N_{m,0} \xredone \dots \xredone N_{m,n}$ এবং $N_{0,n}
  \xredone \dots \xredone N_{m,n}$। লক্ষ করো, $N_{m,0}$ হলো~$P$ এবং
  $N_{0,n}$ হলো~$Q$। $\xred$-এর সংজ্ঞা থেকে উপপাদ্যটি অনুসৃত হয়।
\end{proof}

\end{document}
